% Capa e folha de rosto. No caso do TG, \maketitle gera quatro páginas:
% capa, folha de rosto, ficha catalográfica (CIP) e folha de aprovação.
\maketitle

% Dedicatória
\begin{itadedication}
% PLACEHOLDER: substitua pelo texto da sua dedicatória.
A quem me ensinou que desconfiar de um texto é uma forma de respeito por ele.

\end{itadedication}

% Agradecimentos
\begin{itathanks}
% PLACEHOLDER: substitua pelos seus agradecimentos.
Ao meu orientador, Prof.~Dr.~Cesar Augusto Cavalheiro Marcondes, pela leitura atenta e pela insistência em que o resultado
fosse medido onde ele de fato importa.

Aos professores e colegas do curso, pelas conversas que mudaram o rumo deste
trabalho mais de uma vez.

À minha família, pelo apoio ao longo de toda a graduação.

\end{itathanks}

% Epígrafe
\thispagestyle{empty}
\ifhyperref\pdfbookmark[0]{\nameepigraphe}{epigrafe}\fi
\begin{flushright}
\begin{spacing}{1}
\mbox{}\vfill
% PLACEHOLDER: substitua pela epígrafe que você escolher.
\textit{``Um programa que não distingue o que lhe mandam do que lhe contam
está condenado a obedecer a estranhos.''} \\
\vspace{0.5cm}
\textsc{Autor da epígrafe}

\end{spacing}
\end{flushright}

% Resumo. As palavras-chave ficam aqui, e não em resumo.tex, porque o mesmo
% arquivo alimenta a Folha de Registro do Documento, onde elas não cabem.
\begin{abstract}
\noindent
O uso de LLMs como componentes fixos de software expõe uma superfície de ataque
linguística cuja causa-raiz é a indistinção nativa entre inputs maliciosos e não maliciosos sobre a
entrada do modelo, da qual decorre o \textit{prompt injection}. A defesa,
consolidada cientificamente e comercialmente em implementações industriais,
interpõe um \textit{guardrail} de entrada antes da execução
principal. Esse filtro, porém, é também uma LLM que recebe instrução e dado
pelo mesmo canal, e herda a vulnerabilidade que existe para mitigar. Este
trabalho mede o que se ganha ao aplicar ao próprio filtro a separação entre
canal de instrução e canal de dado que a literatura propõe para o modelo
principal, e ao fazer com que o filtro preserve e marque o trecho suspeito em
vez de removê-lo, repassando a anotação ao componente que executará a tarefa. Essencialmente,
o objetivo é melhorar a segurança de sistemas que utilizem LLM sem recorrer a um modelo \textit{fine-tunado}, mas
adicionando técnicas de segurança que modificam o \textit{input}. O
desenho experimental é $2 \times 2$: delimitação em XML na entrada do
filtro e marcação seletiva na sua saída, totalizando quatro condições, além
de um teste sem filtro como referência inicial e de um filtro
\textit{fine-tunado} da mesma família, avaliado nas mesmas quatro condições,
como ponto de comparação. A avaliação emprega
dois modelos abertos sem \textit{fine-tuning}, cada um instanciando o sistema
inteiro, sobre três \textit{benchmarks} públicos: BIPIA, SEP e NotInject. São
42 execuções por repetição, repetidas dez vezes: 420 execuções e 131.460
registros. O desfecho de cada item é a média das dez repetições, medido na
saída do componente principal, em taxa de sucesso do ataque, utilidade em
tarefas benignas e sobre-bloqueio, com intervalos de confiança por
\textit{bootstrap} pareado sobre os mesmos itens. A combinação dos dois fatores
reduziu a taxa de sucesso do ataque no BIPIA em 0,084 para o Llama-3.1-8B e em
0,139 para o gpt-oss-20b, ambas com intervalo de 95\% afastado do zero.
Sobre o filtro \textit{fine-tunado} a mesma combinação reduziu 0,139, o que
mostra que o ganho não depende de o filtro ser um modelo de uso geral.
Marcar em vez de remover mostrou-se uma
alternativa viável às defesas que removem, com proteção maior que a do modelo
\textit{fine-tunado} da mesma família e com facilidade de implementação muito
maior.

\keywords{\textit{Prompt injection}. \textit{Guardrails} de entrada. Separação
entre instrução e dado. \textit{Spotlighting}. Segurança de LLM. Modelos
abertos.}
\end{abstract}

% Abstract
\begin{englishabstract}
\noindent
Using LLMs as fixed software components exposes a linguistic attack surface
whose root cause is the native inability to distinguish malicious from
non-malicious inputs at the model's input, from which \textit{prompt injection}
follows. The defense, consolidated both scientifically and commercially in
industrial deployments, places an input \textit{guardrail} before the main
execution. That filter, however, is itself an LLM that receives instruction and
data through the same channel, and inherits the very vulnerability it exists to
mitigate. This work measures what is gained by applying to the filter itself
the separation between instruction channel and data channel that the literature
proposes for the main model, and by having the filter preserve and mark the
suspicious span instead of removing it, forwarding the annotation to the
component that will carry out the task. In essence, the goal is to improve the
security of LLM-based systems without resorting to a \textit{fine-tuned} model,
but by adding security techniques that modify the \textit{input}. The
experimental design is $2 \times 2$: XML delimiting at the filter's input and
selective marking at its output, totaling four conditions, plus a no-filter
test as the initial reference and a \textit{fine-tuned} filter from the same
family, evaluated under the same four conditions, as a point of comparison. The
evaluation uses two open models
without \textit{fine-tuning}, each one instantiating the whole system, over
three public \textit{benchmarks}: BIPIA, SEP and NotInject. There are 42 runs
per repetition, repeated ten times: 420 runs and 131,460 records. The outcome
of each item is its mean over the ten repetitions, measured at the output of
the main component, in attack success rate, utility on benign tasks and
over-blocking, with paired \textit{bootstrap} confidence intervals over the
same items. Combining both factors reduced the attack success rate on
BIPIA by 0.084 for Llama-3.1-8B and by 0.139 for gpt-oss-20b, both with a 95\%
interval away from zero. On the \textit{fine-tuned} filter the same
combination reduced it by 0.139, showing that the gain does not depend on the
filter being a general-purpose model. Marking rather than removing proved a
viable alternative to removal-based defenses, with greater protection than the
\textit{fine-tuned} model from the same family and far easier implementation.

\englishkeywords{Prompt injection. Input guardrails. Instruction--data
separation. Spotlighting. LLM security. Open models.}
\end{englishabstract}

% Lista de figuras
\listoffigures

% Lista de tabelas
\listoftables

% Lista de abreviaturas e siglas
\listofabbreviations
\begin{longtable}{ll}
API & \textit{application programming interface} \\
ASR & \textit{attack success rate}, taxa de sucesso do ataque \\
BIPIA & \textit{Benchmark for Indirect Prompt Injection Attacks} \\
CIP & catalogação na publicação \\
FN & falso negativo \\
FP & falso positivo \\
FPR & \textit{false positive rate}, taxa de falsos positivos \\
IC & intervalo de confiança \\
IPI & \textit{indirect prompt injection}, injeção indireta de \textit{prompt} \\
JSON & \textit{JavaScript Object Notation} \\
LLM & \textit{large language model}, modelo de linguagem de grande porte \\
LLMC & \textit{LLM component}, LLM empregada como componente de software \\
MoE & \textit{mixture of experts}, mistura de especialistas \\
RAG & \textit{retrieval-augmented generation} \\
RLHF & \textit{reinforcement learning from human feedback} \\
SEP & \textit{Should It Be Executed or Processed?} \\
TPR & \textit{true positive rate}, taxa de verdadeiros positivos \\
XML & \textit{eXtensible Markup Language} \\
\end{longtable}


% Lista de símbolos
\listofsymbols
\begin{longtable}{ll}
\texttt{\^{}} & marca aposta pelo filtro sobre trecho julgado perigoso \\
\texttt{<data>} & tag que delimita o conteúdo não confiável \\
$\Delta$ & diferença de uma métrica entre uma condição e sua referência \\
$n$ & número de itens executados em uma configuração \\
$p$ & taxa estimada sobre os itens de uma configuração \\
\end{longtable}


% Sumário
\tableofcontents
